\section{Tóm tắt lý thuyết, giải pháp, thuật toán}

\subsection{Tiền xử lý dữ liệu và Phân chia thử nghiệm}

Quá trình tiền xử lý dữ liệu cảm biến được thực hiện chi tiết và riêng biệt cho từng tập dữ liệu, tuân thủ nghiêm ngặt nguyên tắc xử lý tín hiệu số (DSP), chuẩn hóa kích thước cửa sổ và phân chia dữ liệu độc lập theo người dùng (Subject-Independent Split) nhằm loại bỏ rò rỉ dữ liệu.

\subsubsection{Tiền xử lý trên tập dữ liệu MotionSense}
\begin{itemize}
    \item \textbf{Đặc trưng cảm biến trích xuất ($6$ kênh)}: 
    Tín hiệu được thu thập từ điện thoại iPhone 6s đặt tại túi quần trước. $6$ kênh cảm biến bao gồm:
    \begin{enumerate}
        \item $3$ kênh gia tốc tổng ($a_x, a_y, a_z$): Được tính toán bằng cách cộng gia tốc người dùng (\texttt{userAcceleration}) với thành phần trọng lực tĩnh (\texttt{gravity}): 
        $\text{total\_acc} = \text{userAcceleration} + \text{gravity}$ (đơn vị $g \approx 9.8\text{ m/s}^2$).
        \item $3$ kênh vận tốc góc ($\omega_x, \omega_y, \omega_z$): Trích xuất từ cảm biến con quay hồi chuyển (\texttt{rotationRate.x}, \texttt{rotationRate.y}, \texttt{rotationRate.z}, đơn vị rad/s).
    \end{enumerate}
    \item \textbf{Phân đoạn cửa sổ trượt (Sliding Window)}: Tần số lấy mẫu thô của thiết bị là $50\text{ Hz}$. Chuỗi thời gian liên tục được cắt thành các cửa sổ độ dài $L = 128$ mẫu (tương đương $2.56$ giây) với bước trượt $\text{Stride} = 64$ mẫu (độ chồng lấp Overlap $50\%$). Mỗi mẫu dữ liệu sau cắt có dạng Tensor 2D kích thước $(6, 128)$.
    \item \textbf{Phân chia dữ liệu (3-Way Subject-Independent Split)}: Toàn bộ $24$ người tham gia thực nghiệm được chia thành 3 tập hoàn toàn không chồng lấp về đối tượng:
    \begin{itemize}
        \item Tập Train ($14$ subjects: \texttt{sub\_1} đến \texttt{sub\_14}): Thu được $12.375$ cửa sổ.
        \item Tập Validation ($4$ subjects: \texttt{sub\_15} đến \texttt{sub\_18}): Thu được $3.711$ cửa sổ.
        \item Tập Test ($6$ subjects: \texttt{sub\_19} đến \texttt{sub\_24}): Thu được $5.453$ cửa sổ.
    \end{itemize}
    \item \textbf{Ánh xạ nhãn ($6$ lớp)}: Walking ($0$), Upstairs ($1$), Downstairs ($2$), Sitting ($3$), Standing ($4$), Jogging ($5$).
\end{itemize}

\subsubsection{Tiền xử lý trên tập dữ liệu UCI-HAR}
\begin{itemize}
    \item \textbf{Đặc trưng cảm biến trích xuất ($6$ kênh)}: 
    Dữ liệu thu từ thiết bị Samsung Galaxy S II đeo tại vị trí thắt lưng. Trích xuất trực tiếp từ 6 file văn bản thô trong thư mục \texttt{Inertial Signals/}:
    \begin{enumerate}
        \item $3$ kênh gia tốc cơ thể tổng (\texttt{total\_acc\_x}, \texttt{total\_acc\_y}, \texttt{total\_acc\_z}, đơn vị $g$).
        \item $3$ kênh con quay hồi chuyển cơ thể (\texttt{body\_gyro\_x}, \texttt{body\_gyro\_y}, \texttt{body\_gyro\_z}, đơn vị rad/s).
    \end{enumerate}
    \item \textbf{Đóng gói cửa sổ}: Tần số lấy mẫu $50\text{ Hz}$. Dữ liệu thô đã được cắt sẵn ở độ dài $128$ mẫu ($\text{Stride} = 64$). Hệ thống thực hiện ghép (stacking) $6$ chuỗi kênh thành Tensor có kích thước chuẩn $(6, 128)$.
    \item \textbf{Phân chia dữ liệu (Subject-Independent Split)}: Phân chia $30$ người dùng thành các tập độc lập:
    \begin{itemize}
        \item Tập Train ($17$ subjects): Thu được $5.945$ cửa sổ.
        \item Tập Validation ($4$ subjects tách từ tập Train gốc: \texttt{VAL\_SUBJECTS} = $[27, 28, 29, 30]$): Thu được $1.474$ cửa sổ.
        \item Tập Test ($9$ subjects độc lập gốc): Thu được $2.410$ cửa sổ.
    \end{itemize}
    \item \textbf{Ánh xạ nhãn ($6$ lớp)}: Walking ($0$), Upstairs ($1$), Downstairs ($2$), Sitting ($3$), Standing ($4$), Laying ($5$).
\end{itemize}

\subsubsection{Tiền xử lý trên tập dữ liệu HHAR (Heterogeneous HAR)}
\begin{itemize}
    \item \textbf{Đặc trưng cảm biến trích xuất ($6$ kênh)}: 
    Trích xuất đồng bộ từ cụm Gia tốc kế ($x, y, z$) và Con quay hồi chuyển ($x, y, z$) trên các thiết bị điện thoại (\texttt{Phones\_accelerometer.csv}, \texttt{Phones\_gyroscope.csv}) và đồng hồ đeo tay (\texttt{Watch\_accelerometer.csv}, \texttt{Watch\_gyroscope.csv}).
    \item \textbf{Quy trình xử lý tín hiệu số (DSP Pipeline)}:
    \begin{enumerate}
        \item Lọc loại bỏ các nhãn không thuộc $5$ hoạt động chung (bỏ \texttt{bike} và \texttt{null}).
        \item Gom nhóm dữ liệu theo từng phiên ghi thuần nhất \texttt{(User, Device, gt)} và chỉ cắt lấy khoảng thời gian giao nhau (Overlap time) giữa cảm biến gia tốc và con quay.
        \item Xử lý ngắt khoảng thời gian: Cắt đoạn tín hiệu nếu khoảng cách giữa hai mẫu liên tiếp lớn hơn $\Delta t > 1.0$ giây (\texttt{GAP\_THRESHOLD\_S} = $1.0\text{s}$).
        \item Lọc thông thấp (Low-pass Filter): Áp dụng bộ lọc Butterworth bậc $4$ (\texttt{FILTER\_ORDER} = $4$) với tần số cắt $10\text{ Hz}$ (\texttt{LOWPASS\_CUTOFF} = $10.0\text{ Hz}$) nhằm triệt tiêu nhiễu cao tần do rung động thiết bị.
        \item Nội suy tái lấy mẫu (Resampling): Sử dụng nội suy tuyến tính chuyển đổi toàn bộ chuỗi tín hiệu bất đồng bộ về tần số chuẩn $50\text{ Hz}$ (\texttt{TARGET\_FS} = $50.0\text{ Hz}$).
        \item Loại bỏ các phân đoạn thời gian quá ngắn (số mẫu $< 256$).
        \item Cắt cửa sổ trượt $L = 128$ mẫu, $\text{Stride} = 64$ mẫu (Overlap $50\%$) tạo Tensor $(6, 128)$.
    \end{enumerate}
    \item \textbf{Phân chia dữ liệu (Subject-Independent Split)}: Phân chia $9$ người dùng (mã hóa 'a' đến 'i' tương ứng $0 \dots 8$):
    \begin{itemize}
        \item Tập Train ($5$ subjects: $[0, 1, 3, 4, 7]$ tương ứng người dùng 'a', 'b', 'd', 'e', 'h').
        \item Tập Validation ($2$ subjects: $[5, 6]$ tương ứng người dùng 'f', 'g').
        \item Tập Test ($2$ subjects: $[2, 8]$ tương ứng người dùng 'c', 'i').
    \end{itemize}
    \item Đóng gói riêng biệt thành hai tập miền: \texttt{hhar\_phone} (điện thoại) và \texttt{hhar\_watch} (đồng hồ), đồng thời gộp toàn bộ thành tập \texttt{hhar\_combined} phục vụ pretrain SSL.
\end{itemize}

\subsubsection{Chuẩn hóa từng kênh theo mẫu (Channel-wise Instance Normalization)}
Sau khi đóng gói cửa sổ $(6, 128)$, kỹ thuật chuẩn hóa Instance Normalization được áp dụng đồng nhất trên từng kênh cho cả 3 tập dữ liệu:
\begin{equation}
    \hat{\mathbf{x}}_{c, t} = \frac{\mathbf{x}_{c, t} - \mu_{c}}{\sigma_{c} + \epsilon}, \quad \forall c \in \{1, 2, \dots, 6\}, \; t \in \{1, 2, \dots, 128\}
\end{equation}
trong đó $\mu_{c} = \frac{1}{128} \sum_{t=1}^{128} \mathbf{x}_{c, t}$ và $\sigma_{c} = \sqrt{\frac{1}{128} \sum_{t=1}^{128} (\mathbf{x}_{c, t} - \mu_{c})^2}$ là giá trị trung bình và độ lệch chuẩn được tính riêng cho kênh $c$ trên từng cửa sổ $128$ mẫu, và $\epsilon = 10^{-6}$ bảo đảm ổn định số học. Quá trình này giúp triệt tiêu sự chênh lệch về biên độ tuyệt đối giữa các loại cảm biến và các vị trí đeo khác nhau.

\subsection{Các phương pháp Học tự giám sát (SSL)}

\subsubsection{TS-TCC: Temporal \& Contextual Contrasting}
Mô hình TS-TCC (\texttt{TS\_TCC\_Model}) áp dụng cơ chế học tương phản hai mức:
\begin{enumerate}
    \item \textbf{Temporal Contrasting (TC)}: Sử dụng một khối Transformer autoregressive để dự đoán tương lai $k$ bước từ không gian đặc trưng. Cặp dương tính được xác định giữa ngữ cảnh dự đoán của góc nhìn yếu tại thời điểm $t$ và đặc trưng tương lai của góc nhìn mạnh tại thời điểm $t+k$.
    \item \textbf{Contextual Contrasting (CC)}: Áp dụng hàm mất mát NT-Xent (\texttt{NTXentLoss}) trên vector đại diện toàn cục $h_{\text{weak}}$ và $h_{\text{strong}}$ sau lớp Projection Head. Hàm mất mát NT-Xent giữa cặp mẫu thứ $i$ được định nghĩa:
    \begin{equation}
        \mathcal{L}_{\text{NT-Xent}}^{(i, j)} = -\log \frac{\exp(\text{sim}(z_i, z_j) / \tau)}{\sum_{k=1}^{2B} \mathbb{I}_{[k \neq i]} \exp(\text{sim}(z_i, z_k) / \tau)}
    \end{equation}
    với $\text{sim}(u, v) = \frac{u^T v}{\|u\|_2 \|v\|_2}$ là độ tương đồng Cosine và nhiệt độ $\tau = 0.2$.
\end{enumerate}

\subsubsection{SwAV Prototype Learning (Học cụm nguyên mẫu)}
Thay vì so sánh trực tiếp từng cặp mẫu (Instance-level), phương pháp SwAV trong \texttt{PrototypicalHARModel} biểu diễn không gian đặc trưng thông qua $K=45$ nguyên mẫu có thể học $\mathbf{C} \in \mathbb{R}^{K \times D_p}$ (với $D_p = 64$).

Quy trình gán cụm mềm trực tuyến sử dụng thuật toán Sinkhorn-Knopp (\texttt{sinkhorn\_knopp}):
\begin{equation}
    \max_{\mathbf{Q} \in \mathcal{Q}} \text{Tr}(\mathbf{Q}^T \mathbf{C}^T \mathbf{Z}) + \epsilon H(\mathbf{Q})
\end{equation}
trong đó $\mathbf{Q} \in \mathbb{R}^{K \times B}$ là ma trận gán cụm mềm, $H(\mathbf{Q}) = -\sum_{ik} Q_{ik} \log Q_{ik}$ là độ hỗn loạn Entropy, và $\epsilon = 0.05$. Ma trận $\mathbf{Q}$ được tối ưu lặp $3$ vòng để thỏa mãn điều kiện ràng buộc phân bổ đều (Equipartition Constraint).

Hàm mất mát Swapped Prediction (\texttt{SwAVPrototypeLoss}) dự đoán mã cụm của góc nhìn này từ góc nhìn kia:
\begin{equation}
    \mathcal{L}_{\text{SwAV}} = \frac{1}{2} \left[ \ell(\mathbf{q}^{(w)}, \mathbf{p}^{(s)}) + \ell(\mathbf{q}^{(s)}, \mathbf{p}^{(w)}) \right]
\end{equation}
với $\ell(\mathbf{q}, \mathbf{p}) = -\sum_{k} q_k \log p_k$ là hàm Cross-Entropy giữa phân phối gán cụm Sinkhorn $\mathbf{q}$ và phân phối Softmax $\mathbf{p} = \text{Softmax}(\mathbf{C} z / \tau)$.

\subsubsection{CrossHAR: Masked Signal Modeling \& Contrastive Regularization}
CrossHAR kết hợp hai tác vụ phụ (Pretext Tasks):
\begin{itemize}
    \item \textbf{Masked Signal Modeling (MSM)}: Che một tỷ lệ các phân đoạn thời gian trên từng kênh cảm biến bằng \texttt{MaskGenerator}, sau đó sử dụng Encoder-Decoder để tái tạo lại tín hiệu thô ban đầu thông qua hàm mất mát MSE Loss.
    \item \textbf{Contrastive Regularization}: Tạo lực ép điều hòa không gian đặc trưng giữa các biến thể nhiễu biên độ $x^{-}$ và biến thể hoán vị phân đoạn $x^{+}$ nhằm duy trì tính vững chắc của Encoder trước các biến đổi tín hiệu vật lý.
\end{itemize}

\subsection{Kiến trúc các mạng trích xuất đặc trưng (Backbones)}
Nghiên cứu cài đặt và phân tích 4 kiến trúc Backbone trích xuất đặc trưng chính (\texttt{models/encoders/builder.py}):
\begin{enumerate}
    \item \textbf{Standard 1D-CNN (\texttt{StandardSensorEncoder1D})}: Bao gồm 4 khối tích chập 1D liên tiếp với số kênh tăng dần $6 \rightarrow 32 \rightarrow 64 \rightarrow 96 \rightarrow 128$. Sử dụng Kernel kích thước $7, 5, 5, 3$ kết hợp BatchNorm1d, ReLU và MaxPool1d(2). Số lượng tham số toàn mạng xấp xỉ $\sim 178.694$ tham số ($\sim 0.68\text{ MB}$). Kiến trúc sở hữu quán tính quy nạp (Inductive Bias) mạnh đối với dữ liệu cảm biến chuỗi thời gian cục bộ.
    \item \textbf{TSTCC Encoder (\texttt{TSTCCEncoder})}: Tái hiện chính xác $100\%$ kiến trúc \texttt{base\_Model} từ bài báo gốc TS-TCC (IJCAI 2021). Gồm 3 khối Conv1D với Kernel cố định bằng $8$, Stride $1$, Padding $4$, kết hợp Dropout $0.35$ tại Block 1. Đầu ra của Encoder có dạng Feature Map 3D: $(B, 128, 18)$.
    \item \textbf{CNN-Transformer (\texttt{CNNTransformerEncoder})}: Kết hợp các lớp tích chập 1D đầu để trích xuất đặc trưng tần số cao cục bộ, sau đó đưa qua $2$ lớp Transformer Encoder (\texttt{d\_model}=128, \texttt{nhead}=4, \texttt{dim\_feedforward}=256) để học phụ thuộc chuỗi thời gian khoảng cách xa.
    \item \textbf{ViT-1D (\texttt{ViT1DEncoder})}: Biến thể Vision Transformer 1D chia chuỗi thời gian $128$ mẫu thành các mảng Patch độ dài $16$, nhúng Linear Projection và Positional Encoding trước khi đưa qua các lớp Self-Attention.
\end{enumerate}

\subsection{So sánh và đối chiếu lý thuyết}
Bảng~\ref{tab:ssl_comparison} tóm tắt ưu nhược điểm lý thuyết của 3 trường phái SSL được triển khai trong dự án.

\begin{table}[h!]
\centering
\caption{So sánh lý thuyết giữa các hướng tiếp cận SSL cho cảm biến IMU.}
\label{tab:ssl_comparison}
\resizebox{\textwidth}{!}{%
\begin{tabular}{|l|l|l|l|}
\hline
\textbf{Tiêu chí} & \textbf{Instance-level (TS-TCC)} & \textbf{Cluster-level (SwAV Prototype)} & \textbf{Masked Modeling (CrossHAR)} \\ \hline
\textbf{Cơ chế cốt lõi} & So sánh từng cặp mẫu $z_i, z_j$ qua NT-Xent & Gán cụm mềm vào $K=45$ Prototypes & Tái tạo đoạn tín hiệu bị che (MSE) \\ \hline
\textbf{Độ phức tạp tính toán} & $\mathcal{O}(B^2)$ theo kích thước Batch & $\mathcal{O}(B \cdot K)$ với Sinkhorn-Knopp & $\mathcal{O}(B \cdot L)$ tuyến tính theo độ dài chuỗi \\ \hline
\textbf{Hiện tượng rủi ro} & Phụ thuộc lớn vào Kích thước Batch $B$ & Nhạy cảm với số cụm chọn trước $K$ & Có thể học chi tiết nhiễu tần số cao \\ \hline
\textbf{Khả năng đóng góp} & Học đặc trưng toàn cục vĩ mô xuất sắc & Gom nhóm ngữ cảnh tĩnh (Sitting/Standing) & Bảo toàn cấu trúc sóng và nhịp sinh học \\ \hline
\end{tabular}%
}
\end{table}
